\newcommand{\DoNotLoadEpstopdf}{}
\documentclass[11pt,reqno]{amsart}
\usepackage[T1]{fontenc}
\usepackage{lmodern,amsmath,amssymb,amsthm,mathtools,mathrsfs,microtype}
\usepackage[margin=1.1in]{geometry}
\usepackage[hidelinks]{hyperref}
\hypersetup{pdftitle={Collision records, arithmetic, and raw local inversion in a triangular Lorentz gas},pdfauthor={Qian Qi}}
\setlength{\emergencystretch}{2em}
\allowdisplaybreaks[2]
\numberwithin{equation}{section}
\newtheorem{theorem}{Theorem}[section]
\newtheorem{lemma}[theorem]{Lemma}
\newtheorem{proposition}[theorem]{Proposition}
\newtheorem{corollary}[theorem]{Corollary}
\theoremstyle{definition}
\newtheorem{definition}[theorem]{Definition}
\newtheorem{inputstatement}[theorem]{Input statement}
\theoremstyle{remark}\newtheorem{remark}[theorem]{Remark}
\newcommand{\R}{\mathbb R}\newcommand{\Z}{\mathbb Z}
\newcommand{\one}{\mathbf1}\newcommand{\dd}{\,\mathrm d}
\newcommand{\esssup}{\operatorname*{ess\,sup}}
\newcommand{\diver}{\operatorname{div}}\newcommand{\TV}{\operatorname{TV}}
\title[Collision records and raw local inversion]{Collision records, arithmetic, and raw local inversion in a triangular Lorentz gas}
\author{Qian Qi}
\date{October 13, 2026. A2-DYN, revision 95}
\subjclass[2020]{37D50, 37A50, 60F05, 28A75}
\keywords{Lorentz gas, exact collision records, two-ended traces, radial residence, signed pointwise forcing}

\begin{document}
\raggedbottom
\begin{abstract}
We continue the upper pointwise local limit problem for the exact record
of actual returns in a triangular Lorentz gas. On short physical segments,
the original weighted source is approximated in essential height by the
smaller of its two endpoint traces. At every fixed positive entrance
gradient this eliminates the zero-residence contribution of all segments
with a grazing endpoint. A two-sided anchor on the intrinsic radial scale
also estimates complete critical annuli without requiring a prescribed
shell crossing. The residual thin segments cannot originate at a critical
point. Finally we retain signed projection in both roof directions; the
minimum of the two resulting forcings is no larger than the previous
backward forcing and gives an equivalent pointwise criterion. These
results preserve the original source and every exact label. The remaining
two-nongrazing bridge and zero-gradient concentrations are not asserted
to vanish; the unrestricted pointwise theorem remains the objective.
\end{abstract}
\maketitle
\input{journal-problem-v93}

\subsection{Two physical ends rather than one positive inlet}
The target remains \eqref{eq:v76-target}. The new proof addresses the
short-segment and critical-core sources directly. It does not interpolate
a finite norm into an essential supremum or normalize the upper excess.
Theorem~\ref{thm:v95-packets} proves a complete, finite-count comparison
with a two-ended packet. One of its consequences is a vanishing
zero-residence height for every grazing-ended subfamily at fixed positive
entrance gradient. No positive lower residence cutoff is inserted.

Theorem~\ref{thm:v95-anchor} treats targets which remain inside a
previously prescribed critical shell but have enough residence on the
scale of their own gradient. It chooses an anchor on either physical
side and sums the physical energies with multiplicity one. The leading
bound in one annulus is $Cq^{129/128}+C(1+\varepsilon^{-1})q^{257/128}$.
Its finite-count remainder is displayed, so the difference between a
finite annulus sum and a zero-gradient limit is explicit.

Theorem~\ref{thm:v95-bidirectional} retains entrance and terminal signs
through projection in both directions. A critical entrance which was
positive in the backward balance is negative in the forward balance.
Their minimum is no larger than either individual forcing; it is
not a sum of two positive estimates. Its vanishing is equivalent to the
same original scalar law (Theorem~\ref{thm:v95-endpoint}); verification
of that condition still requires control of the remaining concentrations.

\input{journal-inputs-v95}
\input{core/226_two_ended_short_segment_packets}
\input{core/227_intrinsic_radial_residence_anchors}
\input{core/228_bidirectional_signed_pointwise_forcing}

\section{The remaining upper-source estimate}
The short noncritical source is now represented, up to a vanishing
height error, by a packet requiring two nongrazing physical traces.
At fixed positive entrance gradient its grazing-ended part has zero
ordered residence charge. Critical
annuli with sufficient relative residence have an additional
pointwise estimate even when they do not cross the previously
prescribed shell. The complementary relatively thin segments have
strictly positive entrance gradient.

The pointwise law is not concluded from these statements. A family
of positive annular estimates with a common leading constant does
not authorize removal of a lower gradient cutoff inside a collision
limsup. Nor does the minimum of two unknown signed forcings have to
vanish. The complete two-nongrazing short packet and a possible
concentration in disappearing critical annuli remain positive
source-level obligations. All finite-norm, one-sided lower and
conditioning theorems retain their historical scope; unrestricted
point-conditioned laws and arithmetic-zero normalization are not
promoted here. The full preceding article and every inherited proof
remain in the technical archive.

\appendix
\input{appendices/v95_dependency_map}
\input{journal-references-v93}
\end{document}
